\section{The spectrum of a ring}

We shall now associate to every commutative ring a topological
space $\spec R$
in a functorial manner.
That is, there will be a contravariant functor
\[\spec:  \mathbf{CRing} \to \mathbf{Top}  \]
where $\mathbf{Top}$ is the category of topological spaces.
This construction is the basis for scheme-theoretic
algebraic geometry and will be used frequently in the sequel. 

The motivating observation is the following. If $k$ is an algebraically closed
field, then the maximal ideals in $k[x_1, \dots, x_n]$ are of the form
$(x_1-a_1, \dots, x_n  - a_n)$ for $(a_1, \dots, a_n) \in k[x_1, \dots, x_n]$.
This is the Nullstellensatz, which we have not proved yet. We can thus
identify the maximal ideals in the polynomial ring with the space $k^n$.
If $I \subset k[x_1, \dots, x_n]$ is an ideal, then the maximal ideals in
$k[x_1,\dots,x_n]$ correspond to points where everything in $I$ vanishes. See 
\rref{twovarpoly} for a more detailed explanation. Classical affine algebraic
geometry thus studies the set of maximal ideals in an algebra finitely
generated over an algebraically closed field.

The $\mathrm{Spec}$ of a ring is a generalization of this construction.
In general, it is more natural to 
use all prime ideals instead of just maximal ideals.
\subsection{Definition and examples}

We start by defining $\spec$ as a set. We will next
construct the
Zariski topology and later the functoriality.
\begin{definition} 
Let $R$ be a commutative ring. The \textbf{spectrum} of $R$,
denoted $\spec R$, is
the set of prime ideals of $R$.
\end{definition} 

We shall now make $\spec R$ into a topological space. First, we
describe a
collection of sets which will become the closed sets. 
If $I \subset R$ is an ideal, let
\[ V(I) = \left\{\mathfrak{p}: \mathfrak{p} \supset I\right\}
\subset \spec R.
\]

\begin{proposition} 
There is a topology on $\spec R$ such that the closed subsets
are of the form
$V(I)$ for $I \subset R$ an ideal.
\end{proposition} 

\begin{proof} 
Indeed, we have to check the familiar axioms for a topology:
\begin{enumerate}
\item $\emptyset = V((1))$ because no prime contains $1$. So
$\emptyset$ is closed.
\item $\spec R = V((0))$ because any ideal contains zero. So
$\spec R$ is
closed.
\item We show the closed sets are stable under intersections.
Let
$K_{\alpha} = V(I_{\alpha})$ be closed subsets of $\spec R$ for
$\alpha$
ranging over some index set.  Let $I
= \sum I_{\alpha}$. Then 
\[ V(I) = \bigcap K_{\alpha} = \bigcap V(I_{\alpha}),  \]
which follows because $I$ is the smallest ideal containing each
$I_{\alpha}$,
so a prime contains every $I_{\alpha}$ iff it contains $I$.  
\item The union of two closed sets is closed. Indeed, if $K,K'
\subset \spec
R$ are closed, we show $K \cup K'$ is closed. Say $K= V(I), K' =
V(I')$. Then
we claim: 
\[ K \cup K'  = V(II').  \]
Here, as usual, $II'$ is the ideal generated by products $ii', i \in I, i'
\in I'$. If
$\mathfrak{p}$ is \textbf{prime} and contains $II'$, it must
contain one of $I$, $I'$;
this implies the displayed equation above and implies the
result.
\end{enumerate}
\end{proof} 
\begin{definition} 
The topology on $\spec R$ defined above is called the
\textbf{Zariski
topology}. With it,  $\spec R$ is now a topological space.
\end{definition} 

\begin{exercise} 
What is the $\spec$ of the zero ring?
\end{exercise} 

In order to see the geometry of this construction, let us work
several examples.

\begin{example} 
Let $R = \mathbb{Z}$, and consider $\spec \mathbb{Z}$. Then
every prime is generated by one element, since
$\mathbb{Z}$ is a PID. We have that $\spec \mathbb{Z} = \{(0)\}
\cup \bigcup_{p \
\mathrm{prime}} \{ (p)\}$. The picture is that one has all the
familiar primes $(2), (3),
(5), \dots, $ and then a special point $(0)$.

Let us now describe the closed subsets. These are of the form
$V(I)$ where $I
\subset \mathbb{Z}$ is an ideal, so $I = (n)$ for some $n \in
\mathbb{Z}$.

\begin{enumerate}
\item If $n=0$, the closed subset is all of $\spec \mathbb{Z}$.
\item If $n \neq 0$, then $n$ has finitely many prime divisors.
So $V((n))$ consists
of the prime ideals corresponding to these prime divisors.  
\end{enumerate}

The only closed subsets besides the entire space are the finite
subsets
that exclude $(0)$.  
\end{example} 

\begin{example} \label{twovarpoly}
Say $R = \mathbb{C}[x,y]$ is a polynomial ring in two variables.
We will not give a complete description of $\spec R$ here. But we will write
down several
prime ideals.

\begin{enumerate}
\item For every pair of complex numbers $s,t \in \mathbb{C}$,
the collection of polynomials
$f \in R$ such that $f(s,t) = 0$ is a prime ideal $\mathfrak{m}
_{s,t} \subset R$. In
fact, it is maximal, as the residue ring is all of
$\mathbb{C}$. Indeed,
$R/\mathfrak{m}_{s,t} \simeq \mathbb{C}$ under the map $f \to
f(s,t)$.

In fact, 
\begin{theorem}
The $\mathfrak{m}_{s,t}$ are all the maximal ideals in $R$.
\end{theorem} 
This will follow from the \emph{Hilbert Nullstellensatz} to be
proved later
(\rref{gennullstellensatz}).
\item $(0) \subset R$ is a prime ideal since $R$ is a domain. 
\item If $f(x,y) \in R$ is an irreducible polynomial, then $(f)$
is a prime
ideal. This is equivalent to unique factorization in
$R$.\footnote{To be
proved later \rref{}.}  
\end{enumerate}

To draw $\spec R$, we start by drawing $\mathbb{C}^2$, which is identified
with the
collection of
maximal ideals $\mathfrak{m}_{s,t}, s, t \in \mathbb{C}$. $\spec R$ has
additional (non-closed) points
too, as
described above, but for now let us
consider the topology induced on $\mathbb{C}^2$ as a subspace of
$\spec R$.

The closed subsets of $\spec R$ are subsets $V(I)$ where $I$ is
an ideal,
generated by polynomials $\left\{f_{\alpha}(x,y)\right\}$. It is of interest to
determine the subset of $\mathbb{C}^2$ that
$V(I)$
induces. In other words, we ask:
\begin{quote}
What points of $\mathbb{C}^2$ (with $(s,t)$ identified with
$\mathfrak{m}_{s,t}$) lie in $V(I)$?
\end{quote}
Now, by definition, we know that $(s,t)$ corresponds to a point of $V(I)$ if
and only if $I
\subset \mathfrak{m}_{s,t}$.
This is true iff all the
$f_{\alpha} $ lie in $ \mathfrak{m}_{s,t}$, i.e. if
$f_{\alpha}(s,t) =0$ for all
$\alpha$. So the closed subsets of $\mathbb{C}^2$ (with the
induced Zariski
topology) are \emph{precisely the subsets
that can be defined by polynomial equations}.

This is
\textbf{much} coarser
than the usual topology. For instance, $\left\{(z_1,z_2):
\Re(z_1) \geq 0\right\}$ is
not Zariski-closed. 
The Zariski topology is so coarse because one has only algebraic
data (namely,
polynomials, or elements of $R$) to define the topology.
\end{example} 

\begin{exercise} 
Let $R_1, R_2$ be commutative rings. Give $R_1 \times R_2$ a
natural structure
of a ring, and describe $\spec( R_1 \times R_2)$ in terms of
$\spec R_1$ and
$\spec R_2$.
\end{exercise} 


\begin{exercise} 
Let $X$ be a compact Hausdorff space, $C(X)$ the ring of real
continuous
functions $X \to \mathbb{R}$. 
The maximal ideals in $\spec C(X)$ are in bijection with the
points of $X$,
and the topology induced on $X $ (as a subset of $\spec C(X)$ with the Zariski
topology)
is just the usual topology.
\end{exercise}

\begin{exercise}
Prove the following result: if $X, Y$ are compact Hausdorff
spaces and $C(X),
C(Y)$ the associated rings of continuous functions, if $C(X),
C(Y)$ are
isomorphic as $\mathbb{R}$-algebras, then $X$ is homeomorphic to
$Y$.
\end{exercise} 


\subsection{The radical ideal-closed subset correspondence}

We now return to the case of an arbitrary commutative ring $R$.
If $I \subset R$, we get a closed
subset $V(I) \subset \spec R$. It is called $V(I)$ because one
is supposed to
think of it as the places where the elements of $I$ ``vanish,''
as the
elements of $R$ are something like ``functions.'' This analogy
is perhaps best
seen in the example of a polynomial ring over an algebraically
closed field,
e.g. \rref{twovarpoly} above.

The map from ideals into closed sets is very far from being
injective in
general, though by definition it is surjective.

\begin{example} 
If $R = \mathbb{Z}$ and $p$ is prime, then $I = (p), I' = (p^2)$
define the
same subset (namely, $\left\{(p)\right\}$) of
$\spec R$. 
\end{example} 

We now ask why the map from ideals to closed
subsets fails to
be injective. As we shall see, the entire problem disappears if
we restrict to
\emph{radical} ideals.

\begin{definition} 
If $I$ is an ideal, then the \textbf{radical} $\rad(I) $ or $
\sqrt{I}$ is
defined as $$\rad(I) =
\left\{x \in R: x^n \in I \ \mathrm{for} \ \mathrm{some} \ n
\right\}.$$
An ideal is \textbf{radical} if it is equal to its radical.
(This is
equivalent to the earlier \rref{def-radical-ideal}.) 
\end{definition} 

Before proceeding, we must check:
\begin{lemma} 
If $I$ an ideal, so is $\rad(I)$.
\end{lemma} 
\begin{proof} 
Clearly $\rad(I)$ is closed under multiplication since $I$ is.
Suppose $x,y \in \rad(I)$; we show $x+y \in \rad(I)$. Then $x^n,
y^n \in I$
for some $n$ (large) and thus for all larger $n$. The binomial
expansion now
gives
\[ (x+y)^{2n} = x^{2n} + \binom{2n}{1} x^{2n-1}y + \dots +
y^{2n}, \]
where every term contains either $x,y$ with power $ \geq n$, so
every term
belongs to $I$. Thus $(x+y)^{2n} \in I$ and, by definition, we
see then that $x+y \in \rad(I)$.
\end{proof} 

The map $I \to V(I)$ does in fact depend only on the radical of
$I$. In fact, if $I,J$ have the same radical $\rad(I) =
\rad(J)$, then $V(I) = V(J)$.
Indeed, $V(I) = V(\rad(I)) = V(\rad(J)) = V(J)$ by:
\begin{lemma} 
For any $I$, $V(I) = V(\rad(I))$.
\end{lemma} 
\begin{proof} 
Indeed, $I \subset \rad(I)$ and therefore obviously $V(\rad(I))
\subset V(I)$. We have to show the
converse inclusion. Namely, we must prove:
\begin{quote}
If $\mathfrak{p} \supset I$, then $\mathfrak{p} \supset
\rad(I).$
\end{quote}
So suppose $\mathfrak{p} \supset I$ is prime and $x \in
\rad(I)$; then $x^n \in I \subset \mathfrak{p}$ for some $n$.
But $\mathfrak{p}$ is prime, so whenever a product of things
belongs to
$\mathfrak{p}$, a factor does. Thus since $x^n = x \cdot x
\cdots x$, we must
have $x \in \mathfrak{p}$. So
\[ \rad(I) \subset \mathfrak{p},  \]
proving the quoted claim, and thus the lemma.
\end{proof} 

There is a converse to this remark:
\begin{proposition} 
If $V(I) = V(J)$, then $\rad(I) = \rad(J)$. 
\end{proposition} 
So two ideals define the same closed subset iff they have the
same radical.
\begin{proof} 
We write down a formula for $\rad(I)$ that will imply this at
once.
\begin{lemma} \label{radprimescontaining} For a commutative ring $R$ and an
ideal $I \subset
R$,
\[ \rad(I) = \bigcap_{\mathfrak{p} \supset I} \mathfrak{p}.  \]
\end{lemma} 
From this, it follows that $V(I)$ determines $\rad(I)$. This
will thus imply
the proposition.  
We now prove the lemma:
\begin{proof} 
\begin{enumerate}
\item We show $\rad(I) \subset \bigcap_{\mathfrak{p} \in V(I)}
\mathfrak{p} $. In
particular, this follows if we show that if a prime contains
$I$, it contains $\rad(I)$; but we have already
discussed this above.  
\item If $x \notin \rad(I)$, we will show that there is a prime
ideal $\mathfrak{p}
\supset I$ not containing $x$. This will imply the reverse
inclusion and the
lemma.  
\end{enumerate}


We want to find $\mathfrak{p}$ not containing $x$, more
generally not
containing any power of $x$. In particular, we want
$\mathfrak{p} \cap \left\{1,
x, x^2 \dots, \right\} = \emptyset$. This set $S = \left\{1, x,
\dots\right\}$
is multiplicatively closed, in that it contains 1 and is closed
under
finite products. Right now, it does not interset $I$; we want to find
a
\emph{prime} containing $I$ that still does not intersect $\left\{x^n, n
\geq 0\right\}$.


More generally, we will prove:

\begin{sublemma}\label{sublemmamultclosed}
Let $S$ be multiplicatively closed set in any ring $R$ and let
$I$ be any ideal with $I \cap S =
\emptyset$. There is a prime ideal $\mathfrak{p} \supset I$ and
does not
intersect $S$ (in fact, any ideal maximal with respect to the condition of
not intersecting $S$ will do).  
\end{sublemma}
In English, any ideal missing $S$ can be enlarged to a prime
ideal missing $S$.
This is actually fancier version of a previous argument. We
showed earlier that any ideal not
containing the multiplicatively closed subset $\left\{1\right\}$
can be
contained in a prime ideal not containing $1$, in
\rref{anycontainedinmaximal}.

Note that the sublemma clearly implies the lemma when applied to
$S =
\left\{1, x, \dots\right\}.$

\begin{proof}[Proof of the sublemma]
Let $P = \left\{J: J \supset I, J \cap S = \emptyset \right\}$.
Then $P$ is a
poset with respect to inclusion. Note that $P \neq \emptyset$
because $I \in P$. Also,
for any nonempty linearly ordered subset of $P$, the union is in
$P$ (i.e. there is an
upper bound).  
We can invoke Zorn's lemma to get a maximal element of $P$. This
element is an
ideal $\mathfrak{p} \supset I$ with $\mathfrak{p} \cap S =
\emptyset$. We claim
that $\mathfrak{p}$ is prime.

First of all, $1 \notin \mathfrak{p}$ because $1 \in S$. We need
only check
that if $xy \in \mathfrak{p}$, then $x \in \mathfrak{p}$ or $y
\in
\mathfrak{p}$. Suppose otherwise, so $x,y \notin \mathfrak{p}$.
Then $(x,\mathfrak{p}) \notin P$ or
$\mathfrak{p}$ would not be maximal. Ditto for $(y,
\mathfrak{p})$.

In particular, we have that these bigger ideals both intersect
$S$. This means
that there are 
\[ a \in \mathfrak{p} , r \in R \quad \text{such that}\quad a+rx
\in S \]
and 
\[ b \in \mathfrak{p} , r' \in R \quad \text{such that}\quad
b+r'y \in S .\]
Now $S$ is multiplicatively closed, so multiply $(a+rx)(b+r'y)
\in S$.
We find:
\[ ab + ar'y+brx+rr'xy \in S.  \]
Now $a,b \in \mathfrak{p}$ and $xy \in \mathfrak{p}$, so all the
terms above are in $\mathfrak{p}$, and the sum is too. But this contradicts
$\mathfrak{p}
\cap S = \emptyset$. 
\end{proof}
\end{proof} 
\end{proof} 
The upshot of the previous lemmata is:
\begin{proposition}
There is a bijection between the closed subsets of $\spec R$ and
radical ideals
$I \subset R$.
\end{proposition}

\subsection{A meta-observation about prime ideals}

We saw in the previous subsection (\cref{sublemmamultclosed})
that an ideal maximal with respect to the property of not intersecting a
multiplicatively closed subset is prime.
It turns out that this is the case for many such properties of ideals.
A general method of seeing this was developed in \cite{LaRe08}.
In this (optional) subsection, we digress to explain this phenomenon.

If $I$ is an ideal and $a \in R$, we define the notation
\[ (I:a) = \left\{ x\in R: xa \in I\right\} . \]
More generally, if $J$ is an ideal, we define
\[ (I:J) = \left\{x \in R: xJ \subset I\right\} . \]

Let $R$ be a ring, and $\mathcal{F}$ a collection of ideals of $R$. 
We are interested in conditions that will guarantee that the maximal elements
of $\mathcal{F}$ are \emph{prime}.
Actually, we will do the opposite: the following condition will guarantee that
the ideals maximal at \emph{not} being in $\mathcal{F}$ are prime. 

\begin{definition} \label{okafamily}
The family $\mathcal{F}$ is called an \textbf{Oka family} if $R \in
\mathcal{F}$ (where $R$ is considered as an ideal) and whenever $I \subset R$ is an
ideal and $(I:a), (I,a) \in \mathcal{F}$ (for some $a \in R$), then $I \in
\mathcal{F}$.
\end{definition} 

\begin{example} \label{exm:okacard}
Let us begin with a simple observation. If $(I:a)$ is generated by
$a_1, \dots, a_n$ and $(I,a)$ is generated by $a, b_1, \dots, b_m$ (where we
may take
$b_1, \dots, b_m \in I$, without loss of generality), then $I$
is generated by $aa_1, \dots, aa_n, b_1, \dots, b_m$.
To see this, note that if $x \in I$, then $x \in (I,a)$ is a linear
combination of the $\left\{a, b_1, \dots, b_m\right\}$, but the coefficient of
$a$ must
lie in $(I:a)$.

As a result, we may deduce that
the family of finitely generated ideals is an Oka family.
\end{example} 

\begin{example} 
Let us now show that the family of \emph{principal} ideals is an Oka family.
Indeed, suppose $I \subset R$ is an ideal, and $(I,a)$ and $(I:a)$ are
principal.
One can easily check that
$(I:a) = (I: (I, a))$.
Setting $J = (I,a)$, we find that $J$ is principal and $(I:J)$ is too.
However, for \emph{any} principal ideal $J$, and for any ideal $I \subset J$,
\[ I = J (I: J)  \]
as one easily checks. Thus we find in our situation that since $J=(I,a)$ and
$(I:J)$
are principal, $I$ is principal.
\end{example} 

\begin{proposition}[\cite{LaRe08}]\label{okathm} If $\mathcal{F}$ is an Oka
family of
ideals, then any maximal element of the complement of $\mathcal{F}$ is prime.
\end{proposition} 
\begin{proof} 
Suppose $I \notin \mathcal{F}$ is maximal with respect 
to not being in $\mathcal{F}$
but $I$ is  not prime. Note that $I \neq R$ by hypothesis.
Then there is $a \in R$ such that $(I:a), (I,a)$ both strictly contain $I$,
so they must belong to $\mathcal{F}$.
Indeed, we can find $a,b \in R - I$ with $ab \in I$; it follows that $(I,a)
\neq I$ and $(I:a)$ contains $b \notin I$.

By the Oka condition, we have $I \in
\mathcal{F}$, a contradiction.
\end{proof} 

\begin{corollary}[Cohen] \label{primenoetherian}
If every prime ideal of $R$ is finitely generated, then every ideal of $R$ is
finitely generated.\footnote{Later we will say that $R$ is \emph{noetherian.}} 
\end{corollary} 

\begin{proof} 
Suppose that there existed ideals $I \subset R$ which were not finitely
generated.
The union of a totally ordered chain $\left\{I_\alpha\right\}$ of ideals that
are not finitely generated is not finitely
generated; indeed, if $I = \bigcup I_\alpha$ were generated by $a_1, \dots,
a_n$, then all the generators would belong to some $I_\alpha $ and would
consequently generate it.

By Zorn's lemma, there is an ideal maximal with respect to being not finitely
generated. However, by \rref{okathm}, this ideal is necessarily
prime (since the family of finitely generated ideals is an Oka family). This contradicts the hypothesis.
\end{proof} 

\begin{corollary} 
If every prime ideal of $R$ is principal, then every ideal of $R$ is principal.
\end{corollary} 
\begin{proof} 
This is proved in the same way.
\end{proof} 

\begin{exercise} 
Suppose every nonzero prime ideal in $R$ contains a non-zerodivisor. Then $R$
is a domain. (Hint: consider the set $S$ of nonzerodivisors, and argue that
any ideal maximal with respect to not intersecting $S$ is prime. Thus, $(0)$
is prime.)
\end{exercise} 


\begin{remark}
\label{remark-cohen-bound-cardinality}
Let $R$ be a ring. Let $\kappa$ be an infinite cardinal.
By applying
\rref{exm:okacard} and
\rref{okathm}
we see that any ideal maximal with respect to the property of not being
generated by $\kappa$ elements is prime. This result is not so
useful because there exists a ring for which every prime ideal
of $R$ can be generated by $\aleph_0$ elements, but some
ideal cannot. Namely, let $k$ be a field, let $T$ be a set whose
cardinality is greater than $\aleph_0$ and let
\[ R = k[\{x_n\}_{n \geq 1}, \{z_{t, n}\}_{t \in T, n \geq 0}]/
(x_n^2, z_{t, n}^2, x_n z_{t, n} - z_{t, n - 1}) \]
This is a local ring with unique prime ideal
$\mathfrak m = (x_n)$. But the ideal $(z_{t, n})$ cannot
be generated by countably many elements.
\end{remark}

\subsection{Functoriality of $\spec$}
 The construction $R \to \spec R$ is functorial in $R$ in a
contravariant sense. That is, if $f: R \to R'$, there is a
continuous map $\spec
R' \to \spec R$. This map sends $\mathfrak{p} \subset R'$ to
$f^{-1}(\mathfrak{p}) \subset R$, which is easily seen to be a
prime ideal
in $R$. Call this map $F: \spec R' \to \spec R$. So far, we have
seen that
$\spec R$ induces a contravariant functor from $\mathbf{Rings}
\to \mathbf{Sets}$.

\begin{exercise} 
A contravariant functor $F: \mathcal{C} \to \mathbf{Sets}$ (for
some category
$\mathcal{C}$) is called \textbf{representable} if it is
naturally isomorphic
to a functor of the form $X \to \hom(X, X_0)$ for some $X_0 \in
\mathcal{C}$,
or equivalently if the induced covariant functor on
$\mathcal{C}^{\mathrm{op}}$ is corepresentable. 

The functor $R \to \spec R $ is not representable. (Hint:
Indeed, a representable
functor must send the initial object into a one-point set.)
\end{exercise} 

Next, we check that the morphisms induced on $\spec$'s from a
ring-homomorphism are in fact \emph{continuous} maps of
topological spaces.

\begin{proposition} 
$\spec $ induces a contravariant functor from $\mathbf{Rings}$
to the category
$\mathbf{Top}$ of topological spaces.
\end{proposition} 
\begin{proof} Let $f : R \to R'$. 
We need to check that this map $ \spec R' \to \spec R$, which we
call $F$, is
continuous.
That is, we must check that $F^{-1}$ sends closed
subsets of $\spec R$ to closed subsets of $\spec R'$.  

More precisely, if $I \subset
R$ and we take the inverse image $F^{-1}(V(I)) \subset \spec
R'$, it is just
the closed set $V(f(I))$. This is best left to the reader, but
here is the justification. If $\mathfrak{p} \in \spec R'$, then
$F(\mathfrak{p}) = f^{-1}(\mathfrak{p})
\supset I$ if and only if $\mathfrak{p} \supset f(I)$. So
$F(\mathfrak{p}) \in
V(I)$ if and only if $\mathfrak{p} \in V(f(I))$.

\end{proof} 



\begin{example} 
Let $R$ be a commutative ring, $I \subset R$ an ideal, $f: R \to
R/I$. There is a map
of topological spaces
\[ F: \spec (R/I) \to \spec R  .\]
This map is a closed embedding whose image is $V(I)$. Most of
this follows because
there is a bijection between ideals of $R$ containing $I$ and
ideals of $R/I$, and this bijection preserves primality.

\begin{exercise} 
Show that this map $\spec R/I \to \spec R$ is indeed a
homeomorphism from $\spec R/I
\to V(I)$.  
\end{exercise} 
\end{example} 


\subsection{A basis for the Zariski topology}
In the previous section, we were talking about the Zariski
topology. If $R$ is a
commutative ring, we recall that $\spec R$ is defined to be the
collection of
prime ideals in $R$. This has a topology where the closed sets
are the sets of
the form
\[ V(I) = \left\{\mathfrak{p} \in \spec R: \mathfrak{p} \supset
I\right\} . \]
There is another way to describe the Zariski topology in terms
of
\emph{open} sets.  

\begin{definition} 
If $f \in R$, we let 
\[ U_f = \left\{\mathfrak{p}: f \notin \mathfrak{p}\right\}  \]
so that $U_f$ is the subset of $\spec R$ consisting of primes
not containing
$f$. This is the complement of $V((f))$, so it is open.  
\end{definition} 

\begin{proposition} 
The sets $U_f$ form a basis for the Zariski topology. 
\end{proposition} 

\begin{proof} 
Suppose $U \subset \spec R$ is open. We claim that $U$ is a
union of basic
open sets $U_f$. 

Now $U = \spec R - V(I)$ for some ideal $I$.  Then
\[ U = \bigcup_{f \in I} U_f  \]
because if an ideal is not in $V(I)$, then it fails to contain
some $f \in I$,
i.e. is in $U_f$ for that $f$. Alternatively, we could take
complements, whence
the above statement becomes
\[ V(I) = \bigcap_{f \in I} V((f))  \]
which is clear.
\end{proof} 

The basic open sets have nice properties.
\begin{enumerate}
\item $U_1 = \spec R$ because prime ideals are not allowed to
contain the
unit element. 
\item $U_0 = \emptyset$ because every prime ideal contains $0$.
\item $U_{fg} = U_f \cap U_g$ because $fg$ lies in a prime ideal
$\mathfrak{p}$ if and only if one of $f,g$ does.
\end{enumerate}

Now let us describe what the Zariski topology has to do with
localization.
Let $R$ be a ring and $f \in R$. Consider $S = \left\{1, f, f^2,
\dots
\right\}$; this is a multiplicatively closed subset. Last week,
we defined
$S^{-1}R$.

\begin{definition} 
For $S$ the powers of $f$, we write $R_f$ or  $R[f^{-1}]$ for the
localization $S^{-1}R$. 
\end{definition} 

There is  a map $\phi: R \to R[f^{-1}]$ and a corresponding map
\[ \spec R[f^{-1}] \to \spec R  \]
sending a prime $\mathfrak{p} \subset R[f^{-1}]$ to
$\phi^{-1}(\mathfrak{p})$.

\begin{proposition} 
This map induces a homeomorphism of $\spec R[f^{-1}]$ onto $U_f
\subset \spec
R$. 
\end{proposition} 
So if one takes a commutative ring and inverts an element, one
just gets an open
subset of $\spec$. This is why it's called localization: one is
restricting to
an open subset on the $\spec $ level when one inverts something.

\begin{proof}
The reader is encouraged to work this proof out for herself.

\begin{enumerate}
\item 
First, we show that $\spec R[f^{-1}] \to \spec R$ lands in
$U_f$. If
$\mathfrak{p} \subset R[f^{-1}]$, then we must show that the
inverse image
$\phi^{-1}(\mathfrak{p})$ can't contain $f$. If otherwise, that
would imply that
$\phi(f) \in \mathfrak{p}$; however, $\phi(f)$ is invertible,
and then
$\mathfrak{p}$ would be $(1)$.  
\item Let's show that the map surjects onto $U_f$. If
$\mathfrak{p} \subset R$ is a prime
ideal not containing $f$, i.e. $\mathfrak{p} \in U_f$. We want
to construct a
corresponding prime in the ring $R[f^{-1}]$ whose inverse image
is $\mathfrak{p}$.

Let $\mathfrak{p}[f^{-1}]$ be the collection of all fractions
\[ \{\frac{x}{f^n}, x \in \mathfrak{p}\} \subset R[f^{-1}],  \]
which is evidently an ideal. Note that whether the numerator is
in
$\mathfrak{p}$ is \textbf{independent} of the
representing fraction $\frac{x}{f^n}$ used.\footnote{Suppose
$\frac{x}{f^n} =
\frac{y}{f^k}$ for $y \in \mathfrak{p}$. Then there is $N$ such
that
$f^N(f^k x - f^n y) = 0 \in \mathfrak{p}$; since $y \in
\mathfrak{p}$ and $f
\notin \mathfrak{p}$, it follows that $x \in \mathfrak{p}$.}
In fact, $\mathfrak{p}[f^{-1}]$ is a prime ideal. Indeed,
suppose
\[  \frac{a}{f^m} \frac{b}{f^n} \in \mathfrak{p}[f^{-1}] .\]
Then $\frac{ab}{f^{m+n}}$ belongs to this ideal, which means $ab
\in
\mathfrak{p}$; so one of $a,b \in \mathfrak{p}$ and one of the
two fractions
$\frac{a}{f^m}, \frac{b}{f^n}$ belongs to
$\mathfrak{p}[f^{-1}]$. Also, $1/1
\notin \mathfrak{p}[f^{-1}]$.

It is clear that the inverse image of $\mathfrak{p}[f^{-1}]$ is
$\mathfrak{p}$,
because the image of $x \in R$ is $x/1$, and this belongs to
$\mathfrak{p}[f^{-1}]$ precisely when $x \in \mathfrak{p}$.
\item The map $\spec R[f^{-1}] \to \spec R$ is injective.
Suppose
$\mathfrak{p}, \mathfrak{p'}$ are prime ideals in the
localization and the
inverse images are the same.  
We must show that $\mathfrak{p} = \mathfrak{p'}$.

Suppose $\frac{x}{f^n} \in \mathfrak{p}$. Then $x/1 \in
\mathfrak{p}$, so $x
\in \phi^{-1}(\mathfrak{p}) = \phi^{-1}(\mathfrak{p}')$. This
means that $x/1
\in \mathfrak{p}'$, so 
$\frac{x}{f^n} \in \mathfrak{p}'$ too. So a fraction that
belongs to
$\mathfrak{p}$ belongs to $\mathfrak{p}'$. By symmetry the two
ideals must be
the same.  
\item We now know that the map $\psi: \spec R[f^{-1}] \to U_f$
is a continuous
bijection. It is left to see that it is a homeomorphism. We will
show that it
is open.  
In particular, we have to show that a basic open set on the left
side is mapped
to an open set on the right side.
If $y/f^n \in R[f^{-1}]$, we have to show that $U_{y/f^n}
\subset \spec
R[f^{-1}]$ has open image under $\psi$. We'll in fact show what
open set it is.

We claim that
\[ \psi(U_{y/f^n}) = U_{fy} \subset \spec R.  \]
To see this, $\mathfrak{p}$ is contained in $U_{f/y^n}$. This
mean that
$\mathfrak{p}$ doesn't contain $y/f^n$. In particular,
$\mathfrak{p}$ doesn't
contain the multiple $yf/1$. So $\psi(\mathfrak{p})$ doesn't
contain $yf$.
This proves the inclusion $\subset$.  

\item 
To complete the proof of the claim, and
the result, we must show that if $\mathfrak{p} \subset \spec
R[f^{-1}]$ and
$\psi(\mathfrak{p}) = \phi^{-1}(\mathfrak{p}) \in U_{fy}$, then
$y/f^n$ doesn't
belong to $\mathfrak{p}$. (This is kosher and dandy because we
have a bijection.) But the hypothesis implies that $fy \notin
\phi^{-1}(\mathfrak{p})$, so $fy/1 \notin \mathfrak{p}$.
Dividing by $f^{n+1}$
implies that
\[ y/f^{n} \notin \mathfrak{p}  \]
and $\mathfrak{p} \in U_{f/y^n}$. 
\end{enumerate}
\end{proof}

If $\spec R$ is a space, and $f$ is thought of as a ``function''
defined on
$\spec R$, the space $U_f$ is to be thought of as the set of
points where $f$
``doesn't vanish'' or ``is invertible.''
Thinking about rings in terms of their spectra is a very useful
idea. We will bring it up when appropriate.  

\begin{remark} 
The construction $R \to R[f^{-1}]$ as discussed above is an
instance of
localization. More generally, we defined $S^{-1}R$ for $S
\subset R$
multiplicatively closed. We can thus define maps
\( \spec S^{-1}R \to \spec R . \)
To understand $S^{-1}R$, it may help to note that
\[ \varinjlim_{f \in S} R[f^{-1}]  \]
which is a direct limit of rings where one inverts more and more elements.

As an example, consider $S = R - \mathfrak{p}$ for a prime
$\mathfrak{p}$, and for
simplicity that $R$ is countable. We can write $S =
S_0 \cup S_1 \cup \dots$, where each $S_k$ is generated by a
finite number of
elements $f_0, \dots, f_k$. Then $R_{\mathfrak{p}} = \varinjlim
S_k^{-1} R$.
So we have
\[ S^{-1}R = \varinjlim_k R[f_0^{-1} , f_1^{-1}, \dots, f_k^{-1}
] = \varinjlim
R[(f_0\dots f_k)^{-1}]. \]
The functions we invert in this construction are precisely those
which do not
contain $\mathfrak{p}$, or where ``the functions don't vanish.''

The geometric idea is
that to construct $\spec S^{-1}R = \spec R_{\mathfrak{p}}$, we
keep cutting out
from $\spec R$ vanishing locuses of various functions that do
not
intersect $\mathfrak{p}$. In the end, you don't restrict to an
open set, but
to an intersection of them.
\end{remark} 
\begin{exercise} \label{semilocal} 
Say that    $R$ is \emph{semi-local} if it has finitely many maximal ideals.
Let $\mathfrak{p}_1$, \dots, $\mathfrak{p}_n\subset R$ be primes. The complement of
the union, $S=R\smallsetminus \bigcup \mathfrak{p}_i$, is closed under
multiplication, so we can
 localize. $R[S^{-1}] = R_S$ is called the \emph{semi-localization}
 \index{semi-localization} of $R$ at the $\mathfrak{p}_i$.
 
The result of semi-localization is always semi-local. To see this, recall that
the ideals
 in $R_S$ are in bijection with ideals in $R$ contained in $\bigcup
 \mathfrak{p}_i$. Now use prime avoidance.
\end{exercise}

\begin{definition}
For a finitely generated $R$-module $M$, define $\mu_R(M)$ to be the smallest
number
   of elements that can generate $M$.
 \end{definition}
This is not the same as the cardinality of a minimal set of generators. For
example, 2
and 3 are a minimal set of generators for $\mathbb{Z}$ over itself, but
$\mu_\mathbb{Z} (\mathbb{Z}) =1$.

 \begin{theorem}
Let $R$ be semi-local with maximal ideals $\mathfrak{m}_1,\dots,
\mathfrak{m}_n$. Let $k_i = R/\mathfrak{m}_i$. Then
   \[
     \mathfrak{m}u_R(M) = \mathfrak{m}ax \{\dim_{k_i} M/\mathfrak{m}_i M\}
   \]
 \end{theorem}
\begin{proof} 
\add{proof}
\end{proof} 

